\documentclass[10pt,a4paper]{article}
\usepackage[T1]{fontenc}
\usepackage[margin=2cm]{geometry}
\usepackage{isabelle,isabellesym}
\usepackage{amsmath,amssymb}
\usepackage{pdfsetup}

\urlstyle{rm}
\isabellestyle{it}

\begin{document}

\title{S\'andor's Subset-Sum Antichain}
\author{Carles Mar\'in}
\maketitle

\begin{abstract}
A finite set of positive integers is a subset-sum antichain if no two distinct nonempty subsets have sums
dividing one another. Erd\H{o}s~\cite{Er98} reports a construction of S\'andor: for $m\ge 1$ the set
$\{2^i+m\,2^m : 0\le i<m\}$ is such an antichain, which gives antichains of size $(1-o(1))\log_2 n$ in
$\{1,\dots,n\}$ (Erd\H{o}s problem 882). We formalize this. A subset sum of $t$ elements is $tK+b$ with
$K=m\,2^m$ and $b<2^m$ a sum of distinct powers of two; a divisibility $q(tK+b)=t'K+b'$ with $q\ge 2$ forces
$t'=qt$ and $b'=qb$, which is impossible because the number of ones in binary of $qb$ is at most that of $q$
times that of $b$, and that of $q$ is less than $q$. The statement is the one of the Formal Conjectures
project~\cite{FC}.
\end{abstract}

\paragraph{Acknowledgement.} The formalization and the text of this document were prepared with the
assistance of generative AI systems, and checked by the author.

\tableofcontents

\parindent 0pt\parskip 0.5ex

\input{session}

\bibliographystyle{abbrv}
\bibliography{root}

\end{document}
